% ============================================================
% Methodology — Paper A (JNCA): verification procedure for the cascade invariant.
% Procedure, variables, metrics, controls. NO concrete instantiation
% values (those live in the Case Study, Sec.~\ref{sec:results}).
% ============================================================
\section{Methodology}
\label{sec:method}

The procedure turns the cascade invariant into detectors, guards against inflated results, and checks the
metrics against the definition of effectiveness (Sec.~\ref{sec:effectiveness}). Concrete values---population,
time buckets, cascade parameters, hyperparameters---are stated once, in the Case Study (Sec.~\ref{sec:results}).

\subsection{Overview of the verification procedure}
\label{sec:m-overview}
The procedure has four steps that mirror the model. First, instantiate the communication graph $G$
(Def.~\ref{def:cg}) on a population, deriving the contacts $E$ and the temporal rhythm $\tau$. Then build
detectors of increasing temporal awareness, as standard learning instruments over graph-position features, and
two strong practical baselines that read the sending side. Evaluation follows under the leak-aware protocol
below, so that any advantage due to leakage or capacity is destroyed by a control; the last step checks
transfer to real topologies.

\subsection{Leak-aware evaluation protocol}
\label{sec:eval}
Graph detectors are easy to inflate, and the protocol below uses four safeguards against that.

A matched design makes legitimate traffic and attack edges share the same
distribution of endpoints and degrees, meaning they differ only by the temporal pattern and the
degree/identity confounder is removed \citep{arp2022dos}. A time-shuffle control permutes the temporal order of
events, which removes the advantage if it comes from dynamics and leaves it intact if it is an artifact of
topology or identity \citep{pendlebury2019tesseract}; an inductive split by organization checks that nothing is memorized about
the accounts seen in training, and an audit of noisy training labels checks that the learned detector does not
depend on clean labels. Operational metrics report sensitivity at a low FPR, with explicit base-rate
arithmetic, rather than relying on headline AUC alone. A \emph{predictivity probe} then asks whether what was established on one
topology carries over to another (real Enron, SNAP graphs). The probe checks two properties: that the
\emph{qualitative hierarchy} of detectors---which class wins, and in which attack regime it breaks down---is
reproduced on each real graph, and that the headline contrast (receiving-side signal vs.\ volume) keeps its
direction there; with the few seeds available per real graph this is a directional check, not a
significance test. We deliberately do not reduce the probe to a single rank-correlation
coefficient: over a handful of detectors such a coefficient is underpowered, and the transferable claim is the
regime-level contrast, not a full ranking. Together these give the working definition of effectiveness of
Sec.~\ref{sec:effectiveness}: a method works when its advantage meets the standard defined there.

\subsection{Detectors and causal controls}
\label{sec:m-detectors}
The design principle is uniform: features encode graph position and timing, never identity. This makes every detector
inductive and applicable to accounts unseen in training.

\smallskip The invariant operates on the contacts graph
$G=(V,E)$ with time divided into $T$ buckets. An attack is a cascade where a patient-zero $s_0$ sends malicious mail
to contacts; each recipient becomes a carrier with probability $p_\mathrm{inf}$ and propagates to depth $H$
(Fig.~\ref{fig:cascade}). A single event $u\!\to\!v$ is locally indistinguishable from legitimate mail; the
signal is the cascade context, since within the wave's window many of $v$'s neighbors receive mail and $u$ was
hit just before. Four detectors of increasing temporal awareness are compared: a 1-hop tabular model (a lower
bound), a hand-crafted cascade context (1-hop plus burst, freshness, and wave density), a static GNN (topology,
no time), and a temporal GNN with node memory.
% ZMIANA: dodano opis treningu w eksperymencie stealth: kazdy poziom K to osobne zadanie,
% detektor uczy sie i jest sprawdzany w jednym reżimie. Powod: bez tej informacji czytelnik
% mogl podejrzewac, ze zly wynik detektora wolumetrycznego bierze sie z rozjazdu miedzy
% daniem treningowym a testowym.
In the stealth experiment each regime is a self-contained
training task: for every stealth level $K$ the scenario (attack plus matched benign) is generated, the split is
victim-disjoint, and every detector is trained and tested \emph{within that regime only}---no detector ever
sees a mixture of regimes, so its collapse or survival is a property of the regime, not of a train--test
regime shift.

\smallskip The two equations below differ in one respect: whether the model may
look forward in time. The static score summarizes a node by its neighbors in the whole aggregated graph, seeing
who is connected to whom but not in what order. The temporal score keeps a running memory updated
in event order and read from the state before the event it scores, so its memory can only use the past.
% ZMIANA (A1, recenzja wewnetrzna r2): zasada przeczynowosci przepisana z kubełkowej
% (<= b-1) na ZDARZENIOWA — wszystkie detektory widza wylacznie zdarzenia o scisle
% wczesniejszym t(e) = (kubelek, pozycja na liscie). Powod: podzial kubełkowy ukrywal
% fale budujaca sie w biezacym kubełku (asymetria faworyzujaca volume pod blastem);
% regula remisow: porzadek tasowanej listy zdarzen, ustalony per przebieg. TGN czyta
% pamiec per mini-batch (<=64 zdarzen w kubełku) — batch jest wspolczesny. Test
% anty-wyciek: tests/test_causality.py (cechy/logity niezmienne po usunieciu zdarzen
% o czasie >= t(e)).
This rule is enforced uniformly across the whole detector panel (A1): every stream-derived
feature---the volumetric counters, the burst/freshness/propagation windows of the hand-crafted context,
and the precursor/momentum/rare-target components of the path-based baseline---is computed from events
with a strictly smaller total timestamp $t(e)=(b,\text{position})$ only, where the position is the fixed
(shuffled) stream order used as the tie rule within a bucket; the temporal GNN's memory is likewise
updated per mini-batch of at most $64$ simultaneous events in that order, each scored from the state
before it, and the GNNs' static node features (in-/out-degree) are computed from the
\emph{training} split only. Earlier drafts of the package computed some of these features on the full
event stream, and a later draft split causality at bucket boundaries; both hid parts of the signal and
the latter made the baselines asymmetric, so all experiments were
re-run under the uniform event-causal rule and the numbers in this manuscript come from the fixed package.
That past-only
reading is what lets the receiving-side wave become a signal, and it is what a time-shuffle destroys.

\smallskip The static GNN builds embeddings from the aggregated graph with a normalized
adjacency $\hat{A}$ and features $X$,
\begin{equation}
Z = \hat{A}\,\sigma\!\big(\hat{A}\,X\,W_1\big)\,W_2,
\qquad
\mathrm{score}_{\mathrm{stat}}(s,v) = g\!\big([\,Z_s,\,Z_v,\,Z_s\!\odot\!Z_v,\,x_{sv}\,]\big),
\label{eq:static}
\end{equation}
whereas the temporal GNN maintains a memory $m_v$ and scores an event from the memory state
immediately before it in event order, which preserves causality and avoids leakage,
\begin{align}
\mathrm{score}_{\mathrm{temp}}(s,v,e) &= g\!\big([\,m_s^{(<t(e))},\,m_v^{(<t(e))},\,m_s^{(<t(e))}\!\odot\!m_v^{(<t(e))},\,x_{sv}\,]\big),
\label{eq:temp-score}\\
m_v^{(k)} &= \mathrm{GRU}\!\Big(\tfrac{1}{d_v^{(k)}}\!\!\sum_{u:\,(u\to v)\in B_k}\!\!\sigma\!\big(W_m\,m_u\big),\; m_v^{(k-1)}\Big),
\label{eq:temp-mem}
\end{align}
where $B_k$ is the $k$-th mini-batch of events in total time order and $d_v^{(k)}$ the number of its
messages received by $v$ (a normalizer for the mean-pooled message).

The shuffle control is a temporal ablation. Since $\mathrm{score}_{\mathrm{stat}}$ is permutation-invariant in
time while $\mathrm{score}_{\mathrm{temp}}$ is not, the drop in $\mathrm{AUC}_{\mathrm{temp}}$ after time
shuffling indicates how much the model relies on temporal order. The ablation is not perfectly selective:
shuffling also breaks the co-occurrence of events within a bucket, so the drop is an upper bound on the
contribution of order alone (consistent with the shuffled temporal model falling slightly below the static
GNN, Tab.~\ref{tab:detectors}). We use ``causal'' in this ablation sense throughout, not in the sense of
causal identification.

\smallskip Beyond the GNN we implement two strong practical methods, both of which read the sending side.
\begin{itemize}[noitemsep,topsep=2pt]
  \item A volumetric baseline inspired by COMPA \citep{egele2013compa,stringhini2015thataint} scores an event
  by the sender's deviation from its own volume profile---how unusual the recipient, the time bucket, and the
  sender's overall activity are (three per-event features, Supplementary Sec.~S6)---with the same
  gradient-boosting learner as the other tabular detectors,
  \item and a path-based baseline that transfers the causal-path idea of Hopper \citep{ho2021hopper}, a
  detector of lateral movement in login graphs, to the email graph: it looks for a causal precursor,
  momentum, and a rare target.
\end{itemize}
We implement their structural-temporal core rather than the full systems, which is why we label them ``volumetric'' and
``path-based''; the exact feature set of every detector is listed in Supplementary Sec.~S6, together with what
the full systems add on top. Hyperparameters of all detectors are fixed uniformly, without intensive tuning,
so that no detector enjoys a tuning advantage; the concrete settings are given in Supplementary Tab.~S1.

\begin{figure}[pos=tbp]\centering
\resizebox{0.92\textwidth}{!}{%
\begin{tikzpicture}[>=Latex, every node/.style={font=\small},
  acc/.style={circle, draw, minimum size=7mm, inner sep=1pt},
  inf/.style={acc, fill=red!18, draw=red!70!black},
  ben/.style={acc, fill=blue!8},
  mal/.style={-{Latex[length=2mm]}, red!70!black, thick}]
  \node[inf] (p) at (0,0) {$s_0$};
  \node[ben] (a) at (2.4,0.9) {$A$};
  \node[ben] (b) at (2.4,-0.9) {$B$};
  \draw[mal] (p)--(a); \draw[mal] (p)--(b);
  \node[ben] (c) at (4.8,1.6) {$C$};
  \node[ben] (d) at (4.8,0.4) {$D$};
  \node[ben] (e) at (4.8,-0.9) {$E$};
  \draw[mal] (a)--(c); \draw[mal] (a)--(d); \draw[mal] (b)--(e);
  \node[ben] (f) at (7.0,1.0) {$F$};
  \draw[mal] (c)--(f);
  \node[gray] at (0,-2.1) {bucket $t_0$};
  \node[gray] at (3.6,-2.1) {$t_1$};
  \node[gray] at (5.9,-2.1) {$t_2$};
  \draw[gray, dashed] (1.2,-1.8) -- (1.2,2.1);
  \draw[gray, dashed] (3.6,-1.8) -- (3.6,2.1);
  \draw[gray, dashed] (5.9,-1.8) -- (5.9,2.1);
  \node[red!70!black, align=center, font=\scriptsize] at (3.6,2.7)
    {edge $A\!\to\!C$ \emph{locally} = ordinary mail from a known contact;\\
     signal = \emph{synchronized burst} of the wave and sender \emph{freshness}};
\end{tikzpicture}}
\caption[Lateral phishing as a temporal cascade]{Lateral phishing as a temporal cascade: a carrier ($s_0$)
infects contacts, and those in turn infect their own, so the wave propagates across the graph. Detection requires \emph{multi-hop,
temporal} propagation context.}
\label{fig:cascade}
\end{figure}

\subsection{Event model and data split}
\label{sec:m-events}
An event is a pair (sender $s$, victim $v$) with a bucket $b$, labeled legitimate or attack, with two-sided noise so
that the classes are not perfectly separable. Legitimate events are messages between known contacts in rhythm,
with a fraction arriving off-rhythm (an irreducible source of false positives). Attacks are the events of a
cascade seeded at a compromised account, a fraction of which mimic the sender's rhythm. We use two split
regimes, victim-wise (disjoint victim sets) and organization-wise (test organizations disjoint from training
ones, an inductive validation). Headline claims are hardened over many seeds; validations on real graphs use
fewer seeds, since variance is lower on a fixed structure, with a Wilcoxon test where significance matters.
Seed counts for each experiment are given in the respective table or figure caption.

\subsection{Metrics, base rate, and threats to validity}
\label{sec:m-metrics}
The headline metric is sensitivity at a fixed low FPR, the operational working point,
$\text{sensitivity}@\mathrm{FPR} = \max\{\mathrm{TPR}(\theta):\mathrm{FPR}(\theta)\le\alpha\}$ (where $\alpha=0.01$ for our primary analysis); AUC is reported
alongside it. Real systems operate at a far stricter FPR ($<\!4$~FP/million, \citep{ho2019detecting}), so the FPR
reported here (an FPR of $10^{-2}$ against $4\cdot10^{-6}$) is a liberal operating point rather than a
deployment threshold; at a realistic prevalence, precision
follows from Bayes' theorem, which we make explicit in the Case Study (Sec.~\ref{sec:res-rigor}). To isolate the
signal we use the temporal-order shuffle as the causal control: if the signal is truly dynamic, the control
must destroy the advantage.

\smallskip\noindent\textbf{Threats to validity.} The internal threats are label leakage, which we counter by
separating features from the hidden infection state and letting the shuffle control isolate the mechanism, and
dependence on the cascade parameters, which we counter by sweeping the infection probability and depth
(Supplementary Tab.~S2). The external threat is that the maliciousness is synthetic, so we validate the mechanism on
real Enron and two SNAP graphs under controlled injections; the synthetic data lets us control the cascade at
the cost of absolute effects. The construct threat is addressed by preferring sensitivity at low FPR to
headline AUC, in line with \citep{arp2022dos,pendlebury2019tesseract}.
